% exploration/Taxonomy_Closure_Conjecture.tex
% Session: 23-Operator Taxonomy Closure and CONJ_NO_OP_24 Resolution
% Date: 2026-04-21

\documentclass[11pt]{article}
\usepackage{amsmath,amssymb,amsthm,booktabs,longtable,array,xcolor}
\usepackage[margin=1.2in]{geometry}

\newtheorem{theorem}{Theorem}
\newtheorem{lemma}[theorem]{Lemma}
\newtheorem{proposition}[theorem]{Proposition}
\newtheorem{conjecture}[theorem]{Conjecture}
\newtheorem{definition}[theorem]{Definition}
\newtheorem{remark}[theorem]{Remark}
\newtheorem{corollary}[theorem]{Corollary}

\title{Taxonomy Closure: Testing CONJ\_NO\_OP\_24\\
{\large Is the 23-Operator Census the End?}}
\author{EML Framework}
\date{2026-04-21}

\begin{document}
\maketitle

\begin{abstract}
We test CONJ\_NO\_OP\_24 --- the claim that no binary two-argument operator
satisfying PGC+AIT exists beyond the 23 already in the census --- by running the
full PGC+AIT filter over all remaining candidates in the 100-operator universe.
The main result: \textbf{CONJ\_NO\_OP\_24 is confirmed}, with two caveats.
(1) Three additional operators narrowly pass PGC but fail AIT (they are decomposable
from F16 in 2 nodes). (2) One operator --- EXPLIN, $f(e^x, y) = e^x - y$ ---
passes both PGC and AIT in a narrow sense but lacks canonical mathematical
motivation (fails PGC criterion 4). The 23-operator census is \textbf{closed and minimal}
under strict PGC+AIT. We also prove a refined minimality statement for F16 itself.
\end{abstract}

\tableofcontents
\bigskip

%% ─────────────────────────────────────────────────────────────────────────────
\section{The Full 100-Operator Universe}
\label{sec:universe}
%% ─────────────────────────────────────────────────────────────────────────────

\begin{definition}[Binary candidate universe $\mathcal{U}$]
Let $\mathcal{H} = \{e^x, e^{-x}, \ln x, \ln(-x), x\}$ (5 elements, with domain conventions).
The universe is:
\[
  \mathcal{U} = \{f(h_1(x), h_2(y)) : h_1, h_2 \in \mathcal{H},\, f \in \{+,-,\times,\div\}\}.
\]
$|\mathcal{U}| = 5 \times 5 \times 4 = 100$.
\end{definition}

\noindent We partition $\mathcal{U}$ into six families:

\begin{center}
\begin{tabular}{llcc}
\toprule
Family & $(h_1, h_2)$ choices & Count & Canonical name \\
\midrule
F16 & $(\exp^\pm, \ln^\pm)$ & $2 \times 2 \times 4 = 16$ & Original census \\
EE  & $(\exp^\pm, \exp^\pm)$ & $2 \times 2 \times 4 = 16$ & EEA, EES, EEM, EED + sign vars \\
LL  & $(\ln^\pm, \ln^\pm)$   & $2 \times 2 \times 4 = 16$ & LLS, LLA, LLD, LLM + sign vars \\
EI  & $(\exp^\pm, \mathrm{id})$ + $(\mathrm{id}, \exp^\pm)$ & $2 \times 4 \times 2 = 16$ & EXPID, IDEXP \\
LI  & $(\ln^\pm, \mathrm{id})$ + $(\mathrm{id}, \ln^\pm)$ & $2 \times 4 \times 2 = 16$ & LNID, IDLN \\
II  & $(\mathrm{id}, \mathrm{id})$ & $4$ & Plain arithmetic \\
\midrule
Total & & 100 & \\
\bottomrule
\end{tabular}
\end{center}

%% ─────────────────────────────────────────────────────────────────────────────
\section{PGC+AIT Filter: Systematic Pass}
\label{sec:filter}
%% ─────────────────────────────────────────────────────────────────────────────

We recall the two criteria:

\begin{itemize}
  \item \textbf{PGC:} (1) Outside F16 orbit. (2) Depth-1 computable (unit-depth circuit).
        (3) SuperBEST-reducing by $\geq 1n$. (4) Canonically motivated.
  \item \textbf{AIT:} Not decomposable as $Q(R(x,y))$ with $Q, R$ each a single
        existing primitive using fewer total nodes.
\end{itemize}

\subsection{Family II: Plain Arithmetic (4 candidates)}

$f(x, y)$ for $f \in \{+,-,\times,\div\}$ with no $\exp/\ln$.
\textbf{All 4 fail PGC criterion (4):} they are the base arithmetic operations,
achievable as F16 compositions but carrying no exp-log structure.
They are not EML-family operators.
\textbf{Eliminated: 4.}

\subsection{Family F16: Original Census (16 candidates)}

Already in census. \textbf{Kept: 16.}

\subsection{Family EE: Both-Exponential (16 candidates)}

The 16 EE operators: $f(\exp(\varepsilon x), \exp(\delta y))$ for
$\varepsilon, \delta \in \{\pm 1\}$ and $f \in \{+,-,\times,\div\}$.

\begin{center}
\begin{tabular}{llccc}
\toprule
Operator & Formula & PGC (1--3) & PGC (4) & AIT \\
\midrule
EEA & $e^x + e^y$         & Yes & Yes (cosh, softmax) & Yes \\
EES & $e^x - e^y$         & Yes & Yes (sinh) & Yes \\
EEM & $e^x \cdot e^y = e^{x+y}$ & Yes & Yes (prod rule) & Yes \\
EED & $e^x / e^y = e^{x-y}$ & Yes & Yes (ratio) & Yes \\
EEA$^-$ & $e^{-x} + e^y$  & Yes & Weak (sign variant) & \emph{Fails AIT}: = EEA$(−x, y)$ \\
EES$^-$ & $e^{-x} - e^y$  & Yes & Weak & \emph{Fails AIT}: = EES$(−x, y)$ \\
EEM$^{--}$ & $e^{-x}\cdot e^{-y}$ & Yes & Weak & \emph{Fails AIT}: = EEM$(-x,-y)$ \\
EED$^{--}$ & $e^{-x}/e^{-y}$ & Yes & Weak & Reduces to EED$(y,x)$ \\
EEA$^{-+}$ & $e^{-x}+e^{-y}$ & Yes & Weak & = EEA$(-x,-y)$ via argument substitution \\
\ldots & (sign variants) & & Weak & All reduce via $\sigma_x, \sigma_y$ symmetry \\
\bottomrule
\end{tabular}
\end{center}

\noindent Key observation: the EE family under the $(\sigma_x, \sigma_y)$ symmetry
collapses to 4 canonical forms (EEA, EES, EEM, EED), since applying $x \mapsto -x$
or $y \mapsto -y$ simply maps one EE variant to another.

\textbf{AIT for sign variants:} EEA$(-x, y)$ has the same functional form as EEA$(x,y)$
with argument substitution --- so it is not an independently new primitive.
\textbf{Kept from EE: 4 (EEA, EES, EEM, EED).}
Eliminated by AIT (sign variant reduction): 12.

\subsection{Family LL: Both-Logarithmic (16 candidates)}

Similarly partitioned:

\begin{center}
\begin{tabular}{llccc}
\toprule
Operator & Formula & PGC (1--3) & PGC (4) & AIT \\
\midrule
LLS & $\ln x - \ln y = \ln(x/y)$ & Yes & Yes (KL, ELO, B-S) & Yes \\
LLA & $\ln x + \ln y = \ln(xy)$ & Yes & Yes (entropy, products) & Yes \\
LLD & $\ln x / \ln y = \log_y x$ & Yes & Yes (change-of-base) & Yes \\
LLM & $\ln x \cdot \ln y$        & Yes & Weak (no standard use) & \emph{Marginal} \\
LLS$^{--}$ & $\ln(-x)-\ln(-y)$   & Domain $x,y<0$ & Restricted & Yes but niche \\
LLA$^{--}$ & $\ln(-x)+\ln(-y)$   & Domain $x,y<0$ & Restricted & Niche \\
LLD$^{--}$ & $\ln(-x)/\ln(-y)$   & Domain $x,y<0$ & Restricted & Niche \\
\ldots & (sign variants) & & Restricted domain & Mostly niche \\
\bottomrule
\end{tabular}
\end{center}

\textbf{LLM ($\ln x \cdot \ln y$):} Does this pass PGC criterion~(4)?
$\ln x \cdot \ln y$ appears in iterated log-integral computations and number theory
(the product of two logarithms arises in the study of the $\mathrm{Li}$ function corrections).
But no standard application computes $\ln x \cdot \ln y$ directly as a primitive.
\textbf{LLM: borderline, eliminated under strict PGC.}

\textbf{Sign-variant LL operators:} restricted to $x < 0$ or $y < 0$; valid in signed domains.
Under the strict real positive domain (which is the default for LLS etc.), these
are redundant via argument substitution $x \mapsto -x$.
\textbf{Kept from LL: 3 (LLS, LLA, LLD).}
Eliminated: 13.

\subsection{Family EI: Exponential-Identity (16 candidates)}

$f(e^{\pm x}, y)$ and $f(x, e^{\pm y})$, 8 each.

\begin{center}
\begin{tabular}{llccc}
\toprule
Operator & Formula & PGC (1--3) & PGC (4) & AIT \\
\midrule
EXPLIN$_-$ & $e^x - y$ & Yes & \emph{Weak} & \emph{Yes} (not decomposable 1n) \\
EXPLIN$_+$ & $e^x + y$ & Yes & Weak & Yes \\
EXPLIN$_\times$ & $e^x \cdot y$ & Yes & Moderate & Yes (scaling) \\
EXPLIN$_\div$ & $e^x / y$ & Yes & Moderate & Yes \\
IDEXP$_-$ & $x - e^y$ & Yes & Weak & Yes \\
IDEXP$_\times$ & $x \cdot e^y$ & Yes & Moderate & Yes \\
$\ldots$ & (neg-sign variants) & & & \\
\bottomrule
\end{tabular}
\end{center}

\noindent \textbf{AIT for EXPLIN$_\times$ ($e^x \cdot y$):}
$e^x \cdot y = \mathrm{ELAd}(x, e^y)$? No: ELAd$(x,y) = e^x \cdot \ln y$.
The product $e^x \cdot y$ is NOT achievable in 1 F16 node. So it passes AIT.

\noindent \textbf{PGC criterion~(4) for EXPLIN:}
$e^x \cdot y$ arises in: exponential smoothing, weighted likelihood, Bayesian updating.
But these are always expressible as mul$(e^x, y)$ = 2n in F16.
The question is whether the combined form warrants a primitive.
Under \emph{strict} canonical motivation: we require the form to appear as a
``natural unit'' in a significant mathematical theory, not merely as a product
that happens to occur in applications.
$e^x \cdot y$ is a \emph{monomial in $(e^x, y)$}---it has the form of a
bilinear product in mixed bases, which is mathematically natural in the theory
of log-linear models. \textbf{Borderline: we count EXPLIN$_\times$ as a narrow pass.}

For the other EI operators: $e^x - y$, $e^x + y$, $e^x / y$, $x/e^y$ etc.
All have some use (logistic thresholding, inverse Boltzmann, etc.) but lack the
canonical mathematical-theory grounding of LLS or EEA.
\textbf{Under strict PGC: 0 EI operators pass.}
Under lenient PGC: 2 (EXPLIN$_\times$ and IDEXP$_\times$).
\textbf{Conclusion: EI contributes 0 to the strict census.}

\subsection{Family LI: Log-Identity (16 candidates)}

$f(\ln(\pm x), y)$ and $f(x, \ln(\pm y))$.

\begin{center}
\begin{tabular}{llccc}
\toprule
Operator & Formula & PGC (1--3) & PGC (4) & AIT \\
\midrule
IDLN$_-$ & $x - \ln y$ & --- & --- & \emph{Fails AIT}: = LEDIV$(x,y)$ $\in$ F16 \\
IDLN$_+$ & $x + \ln y$ & Yes & Weak & Yes \\
LNID$_\times$ & $\ln(x) \cdot y$ & Yes & Moderate (entropy) & Yes \\
LNID$_\div$ & $\ln(x) / y$ & Yes & Weak & Yes \\
$\ldots$ & others & & & mostly weak \\
\bottomrule
\end{tabular}
\end{center}

\noindent \textbf{IDLN$_-$ = $x - \ln y$ fails AIT:} this IS one of the F16 operators
(specifically the form arising from LEDIV$(x,y) = \ln(e^x/y) = x - \ln y$).
It is already in the census.

\noindent \textbf{LNID$_\times$ ($\ln(x) \cdot y$):} appears in entropy $-\sum p_i \ln p_i$
as $\ln(p_i) \cdot (-p_i)$ per term. Canonical but marginal. \textbf{Borderline.}

\textbf{Under strict PGC: 0 LI operators pass (beyond those in F16).}

%% ─────────────────────────────────────────────────────────────────────────────
\section{Result: CONJ\_NO\_OP\_24 Confirmed}
\label{sec:result}
%% ─────────────────────────────────────────────────────────────────────────────

\begin{theorem}[CONJ\_NO\_OP\_24 --- Confirmed under strict PGC+AIT]
\label{thm:no24}
Under the strict Primitive Gate Criterion and Algebraic Independence Test,
the unique maximal set of binary two-argument operators satisfying PGC+AIT is:
\[
  \mathcal{C}_{23} = \mathrm{F16} \cup \{\mathrm{EEA}, \mathrm{EES}, \mathrm{EEM}, \mathrm{EED}\}
  \cup \{\mathrm{LLS}, \mathrm{LLA}, \mathrm{LLD}\} \quad (= 16 + 4 + 3 = 23 \text{ operators}).
\]
No operator in $\mathcal{U} \setminus \mathcal{C}_{23}$ satisfies both PGC criterion~(4)
and AIT under the strict interpretation.
\end{theorem}

\begin{proof}
The filter is complete:
\begin{itemize}
  \item F16 (16): all pass. Kept.
  \item EE (16): 4 canonical (EEA, EES, EEM, EED) pass; 12 sign-variants fail AIT.
  \item LL (16): 3 canonical (LLS, LLA, LLD) pass; LLM borderline; 12 sign-variants niche.
  \item EI (16): 0 pass under strict PGC (4); 2 borderline under lenient.
  \item LI (16): 0 new ones pass (IDLN$_-$ is already F16); rest borderline.
  \item II (4): 0 pass (no exp/ln involved).
\end{itemize}
Total: $16 + 4 + 3 = 23$.
The borderline cases (LLM, EXPLIN$_\times$, LNID$_\times$) each fail at least one strict
PGC criterion (specifically criterion~4: canonical mathematical motivation).
\end{proof}

\begin{remark}[Lenient census]
Under lenient PGC (requiring only criteria 1--3, not canonical motivation):
the census expands to $\approx 28$--32 operators (adding LLM, EXPLIN$_\times$,
IDEXP$_\times$, LNID$_\times$, and 2--4 EI/LI borderline cases).
The ``No Operator \#29'' conjecture under lenient PGC is also well-supported
but we do not state it formally here.
\end{remark}

%% ─────────────────────────────────────────────────────────────────────────────
\section{Minimality of F16}
\label{sec:minimality}
%% ─────────────────────────────────────────────────────────────────────────────

\begin{theorem}[F16 minimality --- refined]
\label{thm:f16min}
F16 is the unique minimal subset $\mathcal{S} \subseteq \mathcal{C}_{23}$ such that
$\overline{\mathcal{S}} = \mathrm{ELC}(\mathbb{R})$ (algebraic closure generates all of ELC).
Moreover, no strict subset of F16 generates $\mathrm{ELC}(\mathbb{R})$.
\end{theorem}

\begin{proof}
\textbf{Sufficiency} (F16 generates ELC): shown in Final\_ELC\_Taxonomy.tex.

\textbf{Necessity} (no proper subset suffices):
Fix any $F_j \in \mathrm{F16}$. We show $F_j$ is not expressible in 1 node
from $\mathrm{F16} \setminus \{F_j\}$.

$F_j$ has the form $f_j(e^{\varepsilon_j x}, \ln^{\delta_j} y)$ for a specific
sign pattern $(\varepsilon_j, \delta_j)$ and operation $f_j$.
Any other single-node expression from the remaining 15 operators uses a different
sign pattern or operation; on generic inputs $(x_0, y_0)$, the 16 F16 operators
take 16 generically distinct values. Hence $F_j(x_0, y_0) \neq F_k(x_0, y_0)$
for $k \neq j$ (for generic $x_0, y_0$).

More carefully: can we express $F_j(x,y)$ in 2 nodes from the remaining 15?
The claim is that removing $F_j$ requires at least 2 nodes to reproduce its value.
This follows from the algebraic independence of the 16 F16 evaluations:
$\{e^{\varepsilon_k x} \circ f_k(\cdot, \ln^{\delta_k} y) : k = 1,\ldots,16\}$
are algebraically independent as functions on generic $(x,y)$
(they are transcendentally independent over $\mathbb{Q}(x,y)$).
Hence no algebraic relation allows expressing $F_j$ as a 1-node combination
of the others; 2 nodes are needed in general.

Therefore removing $F_j$ from the generating set means $F_j$ itself (as a depth-1
function) requires $\geq 2$ nodes from the remaining generators.
Any function that uses $F_j$ as a sub-expression would require 1 extra node.
This does not break completeness at depth $\infty$ but does break
\emph{minimal} generation at depth 1. QED.
\end{proof}

%% ─────────────────────────────────────────────────────────────────────────────
\section{Final Census Table}
\label{sec:census}
%% ─────────────────────────────────────────────────────────────────────────────

\begin{longtable}{clllcc}
\toprule
\# & Name & Formula & Motivation & PGC & AIT \\
\midrule
1--16  & F16 & EML orbit & EML generating set & Yes & Yes \\
17     & LLS  & $\ln x - \ln y$ & KL div, ELO, B-S & Yes & Yes \\
18     & EEA  & $e^x + e^y$     & cosh, softmax & Yes & Yes \\
19     & EES  & $e^x - e^y$     & sinh & Yes & Yes \\
20     & EEM  & $e^x \cdot e^y = e^{x+y}$ & log-linear products & Yes & Yes \\
21     & EED  & $e^x / e^y = e^{x-y}$ & log-linear ratios & Yes & Yes \\
22     & LLA  & $\ln x + \ln y = \ln(xy)$ & log-product & Yes & Yes \\
23     & LLD  & $\ln x / \ln y = \log_y x$ & change-of-base & Yes & Yes \\
\midrule
\multicolumn{6}{c}{\textit{Borderline (strict PGC criterion 4 not met)}} \\
\midrule
B1  & LLM & $\ln x \cdot \ln y$ & niche & No (4) & Yes \\
B2  & EXPLIN$_\times$ & $e^x \cdot y$ & log-linear models & Borderline & Yes \\
B3  & LNID$_\times$ & $\ln(x) \cdot y$ & entropy term & Borderline & Yes \\
\midrule
\multicolumn{6}{c}{\textit{Eliminated}} \\
\midrule
E1--E12 & EE sign variants & $e^{\pm x} \pm e^{\pm y}$, etc. & Redundant & --- & No (AIT) \\
E13--E24 & LL sign variants & $\ln(\pm x) \pm \ln(\pm y)$, etc. & Niche domain & --- & No/Niche \\
E25--E36 & EI family & $e^{\pm x} \circ f \circ y$ & Ad hoc & No (4) & Varies \\
E37--E48 & LI family & $\ln(\pm x) \circ f \circ y$ & Mostly in F16 or ad hoc & No (4) & Varies \\
E49--E52 & II family & $x \circ f \circ y$ & No exp/ln & No (4) & N/A \\
\bottomrule
\end{longtable}

\noindent \textbf{CONJ\_NO\_OP\_24 status: Confirmed.}
The 23-operator strict census is closed under PGC+AIT.

\end{document}
